\mySmallsection{本章内容}
\begin{itemize}
\item 用简单的基于规则的评分器为 LLM 答案打分
\item 计算 LLM 对自身答案的置信度
\item 编写自我修正循环，让 LLM 迭代改进自己的答案
\end{itemize}

上一章介绍了\emph{推理时扩展 }（简称\emph{推理扩展 }）的概念，它无需进一步训练就能提升模型回复的准确率。该章尤其聚焦于\emph{自洽性（self-consistency）}，即由模型生成多个答案，再取其中出现最频繁的作为最终答案。

如图 5.1 所示，本章将在自洽性之外继续探讨推理扩展，并介绍另一种流行且实用的推理扩展技术——\emph{自我修正（self-refinement）}。与生成多个答案再进行选择不同，自我修正侧重于对单个答案进行迭代精炼，以纠正潜在的错误。

\myGraphic{0.92}{images/CH05_F01_Raschka2.png}{图 5.1 本书涵盖主题的模型图。本章继续第 3 阶段，聚焦于无需额外训练即可提升推理能力的推理时技术，并引入自我修正，即由模型迭代地批判并改进自己的答案。}

\section{对模型回复进行评分与迭代改进}

如前所述，推理扩展是用额外计算换取更高的准确率。到目前为止，我们已经介绍了两种推理扩展技术：\emph{思维链提示}和\emph{自洽性}。

图 5.2 展示的第一种方法是思维链提示，它通过修改提示（prompt）来起作用，例如添加短语 \texttt{"}\texttt{Explain} \texttt{step} \texttt{by} \texttt{step.}\texttt{"}，这样可以促使基础模型写出更长的解释，进而提升答案准确率。对于本身不会自然给出类推理式解释的基础模型，这种方法尤其有用。而训练成的推理模型通常无法从这类推理扩展中受益，因为它们本来就会解释自己的答案。

\myGraphic{0.92}{images/CH05_F02_Raschka2.png}{图 5.2 本书涵盖的三种提升推理的推理时方法。前两种方法已在上一章介绍。本章介绍第三种方法——自我修正，即由模型迭代地改进自己的答案。}

图 5.2 中的第二种方法是自洽性，它让模型并行生成多个答案。然后我们从中选出出现最频繁的候选，作为最终答案。

尽管自洽性相当简单，但它往往能大幅提升答案准确率，因此在准确率优先于延迟的 LLM 应用中成为常见选择。近期的例子包括 DeepSeekMath-V2 和 Google 的 Gemini 3 Deep Think 模式（参考文献见附录 A）。自洽性的一个缺点是，要选出最常见的答案，就要求答案是可供比较的简短答案。

本章将实现一种更通用的技术——\emph{自我修正}，其中 LLM 会迭代地改进自己的答案（图 5.2 中的方法 3）。不过在动手之前，我们得先实现若干评分函数来比较和排序不同的答案，如图 5.3 所示。

\myGraphic{0.92}{images/CH05_F03_Raschka2.png}{图 5.3 我们将构建一个简单的基于规则的评分，计算词元概率和对数概率，并把它们用于一种自我修正方法，让模型迭代地改进自己的答案。}

与前面几章一样，加载预训练的 LLM 之后（图 5.3 中的步骤 1），我们会实现一个简单的\emph{基于规则的评分}函数来说明评分的概念（步骤 2）。接着会讲解\emph{词元概率}（步骤 3）和\emph{词元对数概率}（步骤 4）的概念，它们是实现对数概率评分（logprob 评分）方法（步骤 5）所必需的，而该方法将用于我们的自我修正循环（步骤 6）。

在自我修正循环中，我们主要用对数概率评分函数（图 5.3 中的步骤 5）来跟踪进展。这些评分函数也可用于在自洽性中打破平局，或用来选出最佳回复，而不必依赖多数投票。

图 5.3 中通向自我修正的整个概览看起来相当简短直接，但词元概率和对数概率这两个主题却有些复杂。它们也与下一章相关，届时我们会实现强化学习（reinforcement learning）方法来训练 LLM。为此，这里会花一些时间讨论对数概率评分的概念。

\section{加载预训练模型}

与上一章一样，我们先加载本章要一直使用的模型和分词器。

\begin{codeboxt}{代码清单 5.1 加载分词器和基础模型}
import torch
from reasoning_from_scratch.ch02 import get_device
from reasoning_from_scratch.ch03 import (
     load_model_and_tokenizer
)

device = get_device()
device = torch.device("cpu")  #1

model, tokenizer = load_model_and_tokenizer(
    which_model="base",
    device=device,
    use_compile=False
)
\end{codeboxt}
{\noindent\small\textit{\#1 要在 GPU 上运行代码，请删除此行（前提是你的机器支持）。}\par}

上述代码默认在 CPU 上运行，以确保结果与本章所示的结果总体一致。在 CPU 上仍可能出现微小的数值差异（取决于操作系统和机器），但这些差异通常比在不同加速器后端之间观察到的差异更小、更可预测。

5.4 节和 5.5 节会用到涉及很小数值（小数点后位数很多）的计算，这些计算对设备选择和底层实现细节更为敏感。这在实践中不成问题，但此类不一致在初次阅读时容易让人困惑。因此，建议先用 CPU 设备，之后再考虑 MPS 或 CUDA 设备。

为确认模型已正确加载，我们用上一章的温度和 top-\emph{p} 采样代码，在一个 MATH-500 提示上试试。

\begin{codeboxt}{代码清单 5.2 用温度缩放和 top-\emph{p} 采样生成文本}
from reasoning_from_scratch.ch03 import render_prompt
from reasoning_from_scratch.ch04 import (
    generate_text_stream_concat_flex,
    generate_text_top_p_stream_cache
)

raw_prompt = (
    "Half the value of $3x-9$ is $x+37$. "
    "What is the value of $x$?"
)
prompt = render_prompt(raw_prompt)
prompt_cot = prompt + "\n\nExplain step by step."

torch.manual_seed(0)
response_1 = generate_text_stream_concat_flex(
    model, tokenizer, prompt_cot, device,
    max_new_tokens=2048, verbose=True,
    generate_func=generate_text_top_p_stream_cache,
    temperature=0.9,
    top_p=0.9
)
\end{codeboxt}

模型产生如下回复：

\begin{codebox}
 The problem states that half the value of 3x-9 is equal to x+37. We need to find the value of \(x\).

### Step 2: Translate the problem into an equation
... #1

### Final Answer

\[
\boxed{83}
\]
\end{codebox}
{\noindent\small\textit{\#1 为节省篇幅，回复已截断}\par}

由于我们使用了温度缩放和 top-\emph{p} 采样，改变随机种子会产生不同的回复：

\begin{codebox}
torch.manual_seed(3)
response_2 = generate_text_stream_concat_flex(
    model, tokenizer, prompt_cot, device,
    max_new_tokens=2048, verbose=True,
    generate_func=generate_text_top_p_stream_cache,
    temperature=0.9,
    top_p=0.9,
)
\end{codebox}

这一次，模型的回复如下：

\begin{codebox}
 We start with the given equation:
\[
\frac{1}{2} \times (3x - 9) = x + 37
\]
... #1
Final Answer:
\[
\boxed{83}
\]
\end{codebox}
{\noindent\small\textit{\#1 为节省篇幅，回复已截断}\par}

两种情况下的模型最终答案都正确（\texttt{83}）。第二个回复要短得多，我们可以通过打印每个回复的字符数或词元数来确认：

\begin{codebox}
print("Response 1 characters:", len(response_1))
print("Response 1 tokens:", len(tokenizer.encode(response_1)))
print("\nResponse 2 characters:", len(response_2))
print("Response 2 tokens:", len(tokenizer.encode(response_2)))
\end{codebox}

结果如下：

\begin{codebox}
Response 1 characters: 1419
Response 1 tokens: 534

Response 2 characters: 533
Response 2 tokens: 231
\end{codebox}

较短的回复未必更好。如果两个回复得到相同的正确答案，要判断哪个更好并不容易。这往往取决于人类对清晰度、实用性以及中间步骤部分正确性的偏好。

对回复的中间步骤进行评分仍是一个活跃的研究领域（见附录 A），而过程奖励模型这类评估推理本身的方法，在实践中并不总能带来更好的输出。

如果两个回复在质量上不相上下，有一点是确定的：较短的回复成本更低，因为需要生成的词元更少，因此更受青睐。

\section{用基于规则的评分对 LLM 回复打分}

上一节中，LLM 生成了两个正确的回复。本节将开发一个简单的基于规则的评分函数来比较它们，如图 5.4 所示。

\myGraphic{0.92}{images/CH05_F04_Raschka2.png}{图 5.4 本节实现一个基于规则的评分器，对预训练 LLM 生成的不同答案进行排序。}

基于规则的评分函数会为两个 LLM 回复各打一个分数（见图 5.5），使我们能对它们排序并选出更好的那个。这里的“更好”指的是格式和简洁程度，而非正确性。

\myGraphic{0.92}{images/CH05_F05_Raschka2.png}{图 5.5 两个生成的回复得到相同的正确答案，但解释不同。评分器对回复进行评估，并为每个回复打一个分数。}

如图 5.5 所示，实现评估回复的评分函数（评分器）有多种方式。评分器可以是\emph{启发式的}，即一种简单的基于规则的方法；也可以是另一个为答案打分的 LLM，通常称为 \emph{LLM 作为评判者}（LLM-as-a-judge，见附录 F）；还可以依赖内部的\emph{概率分数}或\emph{似然}，我们将在本章后面进行探讨。

本节先从一个启发式的、基于规则的评分器入手，其实现见代码清单 5.3。之所以从这一版开始，不是因为它最精巧，而是因为它提供了一个简单的基线，便于我们后续比较和理解基于概率的评分器。

\begin{codeboxt}{代码清单 5.3 一个简单的基于规则的评分器}
from reasoning_from_scratch.ch03 import extract_final_candidate
import math

def heuristic_score(
    answer,
    prompt=None,  #1
    brevity_bonus=500.0,
    boxed_bonus=2.0,
    extract_bonus=1.0,
    fulltext_bonus=0.0,
):
    score = 0.0

    cand = extract_final_candidate(answer, fallback="none")    #2
    if cand:
        score += boxed_bonus

    else:     #3
        cand = extract_final_candidate(answer, fallback="number_only")
        if cand:
            score += extract_bonus
        else:
            cand = extract_final_candidate(
                answer, fallback="number_then_full"
            )
            if cand:
                score += fulltext_bonus

    score += 1.5 * math.exp(-len(answer) / brevity_bonus)    #4
    return score
\end{codeboxt}
{\noindent\small\textit{\#1 本节中忽略的一个占位参数 \newline \#2 对带有最终 boxed 值的答案给予奖励。 \newline \#3 若答案不含 boxed 值，则给予较弱的奖励。 \newline \#4 加入随文本长度衰减的简洁性奖励。}\par}

\texttt{heuristic\_score} 函数根据答案能够被多干净地提取出来以及答案的长度，为 LLM 答案分配一个数值分数。例如，如果答案带有 \texttt{\textbackslash boxed\{\}} 形式的回复，我们就给予奖励（\texttt{boxed\_bonus}）。否则，如果它至少包含一个数字，就给予较小的奖励（\texttt{extract\_bonus}）。\texttt{brevity\_bonus} 则根据答案的简短程度赋分。

绝对值并不太重要，更重要的是它们的相对大小。这里 \texttt{boxed\_bonus=2.0} 被有意设得大于 \texttt{extract\_bonus=1.0}，这样带有整齐 boxed 的最终答案会比只能提取出一个数字的答案更受青睐。简洁性项最多贡献 1.5 分，因此它主要用来在质量相近的答案之间打破平局，而不会盖过提取奖励。我们还令 \texttt{fulltext\_bonus=0.0}，这样任意冗长的答案就不会得到奖励，除非我们能在其中提取出更明确的最终候选答案。

\begin{myWarning}{注意}
 这里我们是在开发一个评分器，而不是使用第 3 章的验证器，因为我们要假设并不知道问题的真实答案。只有在用现有测试集评估模型时，我们才使用验证器进行评估。\end{myWarning}

\texttt{prompt} 参数是一个占位参数，在函数中并不起实际作用，但当我们在本章后面开发带有可替换评分器插件的自我修正函数时，它会让事情更轻松（代码也更简单）。

尽管 \texttt{brevity\_bonus=500.0} 看起来数值偏高，但它是通过 \texttt{1.5 * math.exp(-len(answer) / brevity\_bonus)} 作为一个指数衰减项来使用的。下面的绘图结果如图 5.6 所示，展示了 \texttt{brevity\_bonus} 对分数的影响。

\begin{codeboxt}{代码清单 5.4 绘制简洁性惩罚曲线}
import matplotlib.pyplot as plt

def plot_brevity_curve(brevity_bonus, max_len=2048):
    lengths = torch.arange(1, max_len)
    scores = 1.5 * torch.exp(-lengths / brevity_bonus)

    plt.figure(figsize=(4, 3))
    plt.plot(lengths, scores)
    plt.xlabel("Text length (number of characters)")
    plt.ylabel("Score contribution")
    plt.tight_layout()
    plt.show()

plot_brevity_curve(500)
\end{codeboxt}

\myGraphic{0.92}{images/CH05_F06_Raschka2.png}{图 5.6 我们的评分器所用的简单基于规则的长度惩罚。较长的解释得到的分数贡献更小。}

对于较短的答案，分数奖励接近 1.5，而超过 1,000 个字符的答案获得的奖励不超过 0.2。几百个字符的回复相当于一个短段落或一份简洁的解题过程，而 1,000 个字符以上通常意味着相当长的、多句的解释。这一项会略微偏好紧凑的答案，但并不会强迫模型把一切都压缩成一行回复。

为简单起见，简洁性奖励按字符数而非词元数来计算，这样就无需把 \texttt{tokenizer} 传给评分函数。按词元数计算也同样合理，甚至可能更可取。

现在，我们可以把启发式评分器应用到 5.2 节中第一个（较长的）回复上：

\begin{codebox}
print(round(heuristic_score(response_1), 3))
\end{codebox}

得到的分数是 \texttt{2.088}。接下来，我们试试第二个（较短的）回复：

\begin{codebox}
print(round(heuristic_score(response_2), 3))
\end{codebox}

计算出的分数是 \texttt{2.517}，这意味着第二个（较短的）回复会是更受青睐的答案。

现在你已经了解了评分方法的基本思路。在本章后面，我们会开发另一个评分器，并把启发式评分器作为自我修正方法的一部分再次使用。

\begin{mySidebar}{练习 5.1：把启发式评分器用作自洽性中的平局裁决器}

扩展上一章的自洽性实现（\texttt{self\_consistency\_vote}，见代码清单 4.17），使其能够处理候选答案之间的平局。例如，当两个或更多答案得到相同的票数时，就对这几个平局候选应用本节的启发式评分器（\texttt{heuristic\_score}），并选出得分最高的那个。提示：无需修改 \texttt{self\_consistency\_vote} 函数本身来实现平局裁决；改为对 \texttt{self\_consistency\_vote} 返回的 \texttt{results} 字典进行平局裁决即可。

\end{mySidebar}

\begin{mySidebar}{练习 5.2：在 best-of-\emph{N} 设置中使用启发式评分器}

修改自洽性实现，使最终答案用启发式评分器而不是多数投票来选出。为每个问题生成 \emph{N }个候选答案（其中 \emph{N = 2 }或更大），用启发式评分器为每个候选打分，并选出得分最高的作为最终预测。（如果我们用评分器而非多数投票，这种方法在文献中被称为 best-of-\emph{N}，与自洽性相对。）

把这种方法应用在 MATH-500 的一个小子集上，并把结果与普通 best-of-\emph{N} 以及练习 5.1 中的自洽性平局裁决分别比较。

\end{mySidebar}

注意，这个评分器包含若干凭直觉选定的参数，这也是我们称其为启发式评分的原因。要优化提取奖励和简洁性奖励的设置，我们可以把这个评分器接入练习 5.1 和 5.2 的自洽性或 best-of-\emph{N} 方法中，并评估哪些设置在 MATH-500 这样的基准数据集上能带来更高的准确率。

一条实用的经验法则是：如果评分器过于频繁地偏好最终结果难以解析的答案，就提高与提取相关的奖励；如果评分器总是倾向于那些冗长啰嗦、又无助于改善最终结果的答案，就加大简洁性压力。另一方面，如果评分器开始偏好过短但不完整的答案，就降低简洁性惩罚，或提高对易提取的最终答案的奖励。思考这些权衡，往往比孤立地纠结确切的数值更有用。

\section{理解词元概率分数}

现在，我们迈出构建基于模型自身置信度的评分器的第一步（图 5.7 中的步骤 5），这里的置信度是模型赋予的一个概率。在每个位置上，模型会把概率质量分配到可能的下一个词元上。如果某个候选答案中的词元持续获得高概率，说明该答案与模型自身的内部偏好更为契合。反过来，如果模型给答案中的许多词元分配了很低的概率，那就表明模型本身并不强烈支持该答案。

\myGraphic{0.92}{images/CH05_F07_Raschka2.png}{图 5.7 本节将从简单的基于规则的评分器转向词元级别的概率。这些概率将构成 logprob 评分方法的基础，我们稍后会在自我修正方法中使用它。}

本节我们先计算某个候选答案的词元概率分数，并用它们来估计模型认为该答案有多大的可能性（图 5.7 中的步骤 3）。这看起来可能像是绕了弯路，但概率和对数概率这两个概念在下一章又会变得重要起来——届时它们会出现在训练目标内部，而不再只是推理时的评分信号。

本节要计算的词元概率分数，在文献中也被称为\emph{下一个词元概率}、\emph{逐词元概率}、\emph{序列似然}，或更宽泛地称为\emph{词元级似然}。它们表示模型赋予每个可能的下一个词元的概率，以整个词表上的归一化（softmax）分布来表示。

\begin{myWarning}{注意}
 \emph{logprob} 一词是 AI 文献中对\emph{对数概率}（log probability）的常见简写。\end{myWarning}

虽然这番关于概率的讨论听起来可能有些复杂，但它依赖的正是我们上一章讲过的、用于文本生成的同一套机制。例如，在第 4 章中你已经看到，模型在每个阶段都会为词表中的每个词元分配一个 logit 分数，并依据这些分数选出下一个词元。这一过程我们在 4.4 节讨论过，图 5.8 再次予以说明。

\myGraphic{0.92}{images/CH05_F08_Raschka2.png}{图 5.8 LLM 如何选出下一个词元。模型把输入文本转换为词元 ID，为词表中的每个词元计算一个分数，并选出分数最高的词元。右侧的图显示了词表一个子集上的 logit 值，其中 Berlin 对应的词元分数最高。}

上一章中，我们查找了与最高分对应的词表项（词表索引 \texttt{19846}，对应图 5.8 中的词元 \texttt{"}\texttt{Berlin}\texttt{"}），以得到下一个生成的词元。这一次，我们出于不同的目的重新审视 logits：不是要生成下一个词元，而是想用这些分数来量化模型对某个特定答案的置信度。换句话说，这里不会依据这些分数生成任何内容，我们只是在查看它们。

例如，考虑我们想要评分的两个候选答案：

\begin{itemize}
\item \texttt{"}\texttt{The capital of Germany is Berlin}\texttt{"}
\item \texttt{"}\texttt{The capital of Germany is Bridge}\texttt{"}
\end{itemize}

我们的目标是量化模型对每个答案的置信度。有一点很重要：模型的置信度并不自动意味着正确。一个模型

可能仅仅因为某个答案很好地契合了它已经学到的模式，就给该答案分配高概率，即便这个答案在事实上是错误的。当这种情况发生时，我们仍然需要其他方法来检测并纠正这种自信却错误的答案，例如检索增强生成（retrieval-augmented generation）、网络搜索、计算器或代码执行等外部工具，或自洽性这类基于比较的方法。

在图 5.8 中，输入文本 \texttt{"}\texttt{The capital of Germany is}\texttt{"} 被送入模型，得分最高的下一个词元（\texttt{"}\texttt{Berlin}\texttt{"}）被选中。这里我们不是要选出一个词元，而是要比较分配给两个候选下一个词元 \texttt{"}\texttt{Berlin}\texttt{"} 和 \texttt{"}\texttt{Bridge}\texttt{"} 的分数，如图 5.9 所示。（之所以选用 \texttt{"}\texttt{Bridge}\texttt{"} 作为对照，是因为它出现在图 5.8 所示的词表范围内。）

\myGraphic{0.92}{images/CH05_F09_Raschka2.png}{图 5.9 我们如何查找特定词元的 logit 分数。把输入文本送入模型后，我们会得到词表中每个词元的 logit 值。接着，我们把想要评分的候选词元转换为词元 ID，并从分布中读出它们对应的 logit 值。}

我们不直接使用图 5.9 中所示的原始 logits，而是把它们转换为概率。概率更易于解读，在不同输入之间可以相互比较，而且它们构成了本章后面要使用的 logprob 评分方法的基础。

如第 4 章所讨论的，应用 \texttt{torch.softmax} 可以把 logits 转换为概率值。图 5.10 显示了得到的概率分布。图 5.10 中所示的词元概率是用 \texttt{torch.softmax} 计算得到的，与我们 4.4.3 节中作为多项式采样过程一部分所做的一样。不同之处在于我们并不从该分布中采样，而只是查找分配给特定词元的概率。

\myGraphic{0.92}{images/CH05_F10_Raschka2.png}{图 5.10 下一个词元评分。输入文本被转换为词元 ID 并送入 LLM，后者输出下一个词元的 logits。经过 softmax 之后，这些 logits 变为概率，其中像 \texttt{"}\texttt{Berlin}\texttt{"}这样的词元获得高概率，而不太可能的替代项如 \texttt{"}\texttt{Bridge}\texttt{"}则得到接近零的值。}

在图 5.10 中，只有 \texttt{"}\texttt{Berlin}\texttt{"} 的条形可见，其概率为 \texttt{0.1695}。绘图范围内其他词元的概率都接近零。这表明，对于输入文本 \texttt{"}\texttt{The capital of Germany is}\texttt{"}，模型赋予 \texttt{"}\texttt{Berlin}\texttt{"} 作为下一个词元的置信度高于 \texttt{"}\texttt{Bridge}\texttt{"}。

\begin{myWarning}{注意}
 图 5.10 中所示的条形并不求和为 1，因为完整词表包含 151,000 个词元，而该图只显示了分布的一小部分（词元索引 19,800–19,900）。\end{myWarning}

在实践中，我们通常希望比较完整的答案，而不是单个词元。图 5.11 说明了这种从单词元评分到序列级评分的扩展：我们在其中计算每个词元在给定其前序词元条件下的概率。除对序列中的每个词元都要重复执行外，这与图 5.10 中的过程完全相同。

\myGraphic{0.92}{images/CH05_F11_Raschka2.png}{图 5.11 计算给定序列的词元概率分数。对于每个位置，我们把前面的文本送入模型，并读出下一个词元的 softmax 概率。把这些条件概率相乘，就得到整个序列的联合概率。}

在普通的文本生成过程中，模型是一边推进一边一步一步地选出下一个词元的。这里，我们以不同的方式使用同一套机制：不做新词元采样，而是取一个已完成的候选答案，在模型下逐个为它的每个词元打分。然后把得到的这些概率相乘，得出整个序列的概率，也称为\emph{联合概率}。

\begin{mySidebar}{词元序列的联合概率}

对于熟悉数学记号的人，词元序列的联合概率可以简洁地写成\emph{条件概率}的乘积。对于序列 𝑥1、𝑥2、...、𝑥𝑇 以及模型权重 \emph{W}，它等于

\end{mySidebar}

\myGraphic{0.7}{images/eq-chapter-5-svg-eq-95.png}{}

\begin{mySidebarBox}

展开后，这变成

\end{mySidebarBox}

\myGraphic{0.7}{images/eq-chapter-5-svg-eq-97.png}{}

\begin{mySidebarBox}

这与我们在图 5.11 中所计算的完全一致。在每个位置，我们把上下文送入模型，得到下一个词元的概率，再把这些分数沿整个序列相乘。

\end{mySidebarBox}

理想情况下，赋予序列 \texttt{"}\texttt{The} \texttt{capital} \texttt{of} \texttt{Germany} \texttt{is} \texttt{Berlin}\texttt{"} 的概率应远高于那个毫无意义的答案 \texttt{"}\texttt{The} \texttt{capital} \texttt{of} \texttt{Germany is} \texttt{Bridge}\texttt{"}。由于联合概率是许多小数值相乘得到的，图 5.11 中两个序列的结果概率都接近于零。我们将在下一节解决这个问题。

本例中的两个答案几乎完全相同，只在最后一个词元上不同（\texttt{"}\texttt{Berlin}\texttt{"} 与 \texttt{"}\texttt{Bridge}\texttt{"}）。同样的方法也可以用于差异更大的序列，例如 5.2 节中针对 MATH-500 提示生成的答案（“\texttt{Half} \texttt{the} \texttt{value} \texttt{of} \texttt{\$3x-9\$ is \$x+37\$. What} \texttt{is} \texttt{the} \texttt{value} \texttt{of} \texttt{\$x\$?}\texttt{"}）。

\begin{mySidebar}{这种概率评分与贪心采样以及温度加 top-\emph{p} 采样有何不同}

当我们计算词元概率分数时，并不是在生成文本。完整序列已经存在，我们只是在给定固定上下文作为输入的情况下，向模型查询每个下一个词元的概率。这些概率不会影响后续的输入。这是一种回顾式的评分过程，而不是生成步骤。

我们在第 4 章编码的生成方法，例如 \texttt{generate\_text\_stream\_cache} 函数（贪心采样）和 \texttt{generate\_text\_top\_p\_stream\_cache} 函数，行为则大不相同。贪心采样总是选出最可能的下一个词元，而温度采样和 top-\emph{p} 采样则从经过重塑的概率分布中抽取。这些过程会在每一步修改分布，然后确定一个单独的词元，该词元又成为下一个位置的输入。

这些采样策略（包括贪心采样）都不能保证产生模型下整体概率最高的序列。在某一步看起来并非最优的词元，可能会导致后续某一步中接下来的词元概率要大得多。反过来，一个局部最优的选择也可能在之后把模型引向低概率的序列。

即便我们真能找到模型下全局概率最高的序列，那也不能保证它就是正确答案，或者对用户最有用的答案。它只意味着，在模型学到的分布下，这是模型最偏好的序列。

温度采样和 top-\emph{p} 采样在抽取样本之前会通过展平或截断分布进一步增加变异性，这可能使生成的文本与全局最可能的序列更加脱节。

通过对一个已有答案进行评分，我们避开了所有这些麻烦。我们不运行采样算法，模型也不选择任何词元。我们只是评估：如果给定的序列已经被写出，它本会有多大的概率。

\end{mySidebar}

我们已经看过了概念性的解释，现在就来看看词元概率计算的实际运行。

\begin{codeboxt}{代码清单 5.5 计算下一个词元的概率}
@torch.inference_mode()
def calc_next_token_probas(model, tokenizer, prompt, device, show=True):

    token_ids = torch.tensor(tokenizer.encode(prompt), device=device)

    logits = model(token_ids.unsqueeze(0)).squeeze(0)     #1
    all_probas = torch.softmax(logits, dim=-1)             #1

    t_idx = torch.arange(0, token_ids.shape[0] - 1, device=device) #2

    next_ids = token_ids[1:] #3

    next_token_probas = all_probas[t_idx, next_ids] #4

    prod_next_token_probas = torch.prod(next_token_probas)  #5

    if show:
        print("Next-token probabilities:", next_token_probas)
        print("Joint probability:", prod_next_token_probas)

    else:
        return next_token_probas, prod_next_token_probas
\end{codeboxt}
{\noindent\small\textit{\#1 像文本生成函数中那样获取 logits 和概率。 \newline \#2 选择要评分的位置（这里：全部）。 \newline \#3 由于我们已有文本，因此知道真正的下一个词元。 \newline \#4 获取每个下一个词元的概率。 \newline \#5 序列的似然是各概率分数之积。}\par}

如你所见，计算下一个词元概率的 \texttt{calc\_next\_token\_probas} 函数（图 5.11 中有说明）相对简短，乍看之下很简单。但为了完成计算，这个函数里发生了很多事情。图 5.12 用一个更简单的输入文本示例和一个只有五个词的小词表对此加以说明。

\myGraphic{0.92}{images/CH05_F12_Raschka2.png}{图 5.12 提取下一个词元的概率。把输入文本转换为词元 ID 后，模型计算 logits，并用 softmax 函数将其转换为概率。接着，利用表示位置和真正的下一个词元的索引张量，我们就能得到模型为每个下一个词元计算出的概率。}

\texttt{calc\_next\_token\_probas} 函数（图 5.12 中有说明）分几步计算下一个词元的概率。首先，它以与前面介绍的文本生成函数相同的方式计算 logits（例如 2.7 节 \texttt{generate\_text\_basic} 中的 \texttt{out = model(token\_ids)} 这一行）。

得到的 logits（图 5.12 中的步骤 2）形状为 \texttt{[sequence\_length, vocab\_size]}，随后用 \texttt{torch.softmax} 转换为归一化的概率分布（步骤 3），使每个位置上的概率之和为 1。

为了提取每个位置上分配给真正下一个词元的概率，我们要构造一个索引张量 \texttt{t\_idx}，用于我们想要评分的位置；再构造另一个张量 \texttt{next\_ids}，用于存放对应的目标词元。例如，如果输入文本是 \texttt{"}\texttt{The} \texttt{capital} \texttt{of} \texttt{Germany} \texttt{is} \texttt{Berlin}\texttt{"}，那么 \texttt{t\_idx} 指向 \texttt{"}\texttt{The} \texttt{capital} \texttt{of} \texttt{Germany} \texttt{is}\texttt{"} 对应的位置，而 \texttt{next\_ids} 则包含向后平移一位的词元：\texttt{"}\texttt{capital} \texttt{of} \texttt{Germany} \texttt{is} \texttt{Berlin}\texttt{"}。

这些目标词元其实就是把输入词元向后平移一位，因为 LLM 的训练目标就是预测序列中的下一个词元。例如，给定输入 \texttt{"}\texttt{The} \texttt{capital} \texttt{of} \texttt{Germany} \texttt{is} \texttt{Berlin}\texttt{"}，模型会在位置 2（\texttt{"}\texttt{capital}\texttt{"} 对应的词元）被要求预测位置 3（\texttt{"}\texttt{of}\texttt{"} 对应的词元）。使用 \texttt{all\_probas[t\_idx,} \texttt{next\_ids]}（步骤 4 和步骤 5）就能准确地取出序列中每个位置上的这些概率值。

最后，函数使用 \texttt{torch.prod}，把逐词元概率的乘积作为序列似然（步骤 6）。

我们来看看这个函数的实际运行：

\begin{codebox}
torch.set_printoptions(precision=4, sci_mode=True)
calc_next_token_probas(
        model, tokenizer, device=device,
        prompt="The capital of Germany is Berlin"
    )
\end{codebox}

得到的输出如下：

\begin{codebox}
Next-token probabilities: tensor([6.1512e-05, 4.6484e-01,
1.6724e-02, 7.3828e-01, 1.6895e-01], dtype=torch.bfloat16)
Joint probability: tensor(5.9372e-08, dtype=torch.bfloat16)
\end{codebox}

我们再试另一段文本：

\begin{codebox}
calc_next_token_probas(
        model, tokenizer, device=device,
        prompt="The capital of Germany is Bridge"
    )
\end{codebox}

结果输出如下：

\begin{codebox}
Next-token probabilities: tensor([6.1512e-05, 4.6484e-01, 1.6724e-02,
7.3828e-01, 2.9802e-07], dtype=torch.bfloat16)
Joint probability: tensor(1.0481e-13, dtype=torch.bfloat16)
\end{codebox}

分析这些结果可见，第一个回复（以 \texttt{"}\texttt{Berlin}\texttt{"} 结尾）的联合概率（序列似然）为 \texttt{5.9372e-08}，大于第二个回复的联合概率 \texttt{1.0481e-13}。（用小数表示，\texttt{5.9372e-08} 即 \texttt{0.000000059372}，\texttt{1.0481e-13} 即 \texttt{0.00000000000010481}。）

这些结果是合理的，因为 \texttt{"}\texttt{Berlin}\texttt{"} 这个答案得到的序列似然高于 \texttt{"}\texttt{Bridge}\texttt{"} 这个答案。两个值都极其小，因为把许多小于 1 的概率相乘，很快就会得到趋近于零的数。这使得原始似然在实践中难以使用，尤其对于更长的序列，下溢会成为一个问题。

\begin{myWarning}{注意}
 这些分数反映的是模型的内部似然，而不是经过校准的正确概率。换句话说，它们告诉我们的是，在模型自身的分布下，它有多偏好一个答案甚于另一个答案，而不是该答案是否真的为真、正确或有用。分数高意味着该答案很好地契合了模型学到的模式，并不代表它一定正确。\end{myWarning}

下一节我们将用对数概率改造这一计算，从而避开这些数值问题，并为整个序列的评分提供一种更稳定的方式。

\begin{mySidebar}{概率与似然}

你可能已经注意到，我们同时使用了“概率”（probability）和“似然”（likelihood）这两个词。它们有区别吗？在统计学中，概率描述的是在我们观察到任何数据之前某事件发生的可能性，并且所有可能结果的概率之和必须为 1。与之相对，似然衡量的是某个特定模型对观测数据的解释有多好，它被视为模型权重或参数的函数，而不是数据的函数。简而言之，概率是在给定模型下预测数据，而似然是在给定数据下评估模型。（更详细的示例，见我这篇关于概率与似然的文章：“What is the difference between likelihood and probability?”，网址为 \href{https://sebastianraschka.com/faq/docs/probability-vs-likelihood.xhtml}{https://sebastianraschka.com/faq/docs/probability-vs-likelihood.xhtml}。）

在 LLM 中，下一个词元概率严格来说是一个概率，而不是似然，因为它来自所有可能下一个词元上的归一化分布，其总和为 1。换句话说，\texttt{next\_token\_probas} 中的值是概率，而不是似然。它们直接来自模型在每个位置上对整个词表做的 softmax，这是一个归一化的概率分布。

以联合概率形式计算出的乘积（\texttt{torch.prod(next\_token\_probas)}）就是模型赋予该序列的概率。这个量通常被称为序列似然（尤其是当把它看作模型参数的函数时）。

\end{mySidebar}

\section{从词元概率分数到对数概率}

上一节计算出的词元概率可以用作评分函数，对不同的回复进行排序。如我们所见，这些概率分数可能非常小，尤其是在相乘得到联合概率或序列似然之后。

本节中，为避免处理这类极小数值时常常出现的数值稳定性问题，我们将对这些概率应用\emph{对数缩放}（也称为\emph{log 缩放}），如图 5.13 的概览所示。

\myGraphic{0.92}{images/CH05_F13_Raschka2.png}{图 5.13 从词元概率转向词元对数概率的概览。后者为自我修正中稍后使用的对数概率评分提供了数值上更稳定的基础。}

\begin{myWarning}{注意}
 本节的重点是对上一节的概率施加一种缩放变换。我们的总体目标仍然是计算词元概率，或者在这个情形下是计算对数概率（你很快就会看到）。\end{myWarning}

考虑下面这个为若干 logit 值计算概率的简单示例：

\begin{codebox}
torch.set_printoptions(precision=4, sci_mode=False)
logits = torch.linspace(-2, 2, steps=7)
probas = torch.softmax(logits, dim=-1)
print(probas)
\end{codebox}

概率值如下：

\begin{codebox}
tensor([0.0090, 0.0175, 0.0341, 0.0665, 0.1295, 0.2522, 0.4912])
\end{codebox}

我们可以通过 \texttt{torch.log} 函数把它们转换为对数概率分数：

\begin{codebox}
print(torch.log(probas))
\end{codebox}

对数概率值如下：

\begin{codebox}
tensor([-4.7109, -4.0442, -3.3776, -2.7109, -2.0442, -1.3776, -0.7109])
\end{codebox}

这里，\texttt{torch.log} 应用的是自然对数，其中 𝑙𝑜𝑔(0.0090) = –4.7109，反之，𝑒–4.7109 = 0.0090。

除了串联使用 \texttt{torch.log(torch.softmax(...))}，PyTorch 还提供了一个优化的 \texttt{torch.log\_softmax} 函数，把这两个操作合并在一起：

\begin{codebox}
log_probas = torch.log_softmax(logits, dim=-1)
print(log_probas)
\end{codebox}

与之前类似，它返回如下结果：

\begin{codebox}
tensor([-4.7109, -4.0442, -3.3776, -2.7109, -2.0442, -1.3776, -0.7109])
\end{codebox}

注意，对数缩放只改变数值的大小，不改变它们的顺序，因此似然更高的序列，其对数似然也一定更高。为了直观地看到这一点，我们把 logits、概率和对数概率并排绘制出来。

\begin{codeboxt}{代码清单 5.6 绘制 logits、softmax 概率和 log-softmax 值}
plt.figure(figsize=(9, 4))

plt.subplot(1, 3, 1)#1
plt.bar(range(len(logits)), logits, color="C0", alpha=0.7) #1
plt.title("Logits") #1
plt.xlabel("Token index") #1
plt.ylabel("Value") #1
plt.grid(alpha=0.3) #1
 #1
plt.subplot(1, 3, 2)#2
plt.bar(range(len(probas)), probas, color="C1", alpha=0.7)
plt.title("torch.softmax(logits)")
plt.xlabel("Token index")
plt.ylabel("Probability")
plt.ylim(0, 1)#3
plt.grid(alpha=0.3)

plt.subplot(1, 3, 3)#4
plt.bar(range(len(log_probas)), log_probas, color="C2", alpha=0.7) #4
plt.title("torch.log_softmax(logits)") #4
plt.xlabel("Token index") #4
plt.ylabel("Log-probability") #4
plt.grid(alpha=0.3) #4
 #4
plt.tight_layout()
plt.savefig("logits_softmax_log_softmax.pdf")
plt.show()
\end{codeboxt}
{\noindent\small\textit{\#1 绘制 logits \newline \#2 绘制 softmax 概率 \newline \#3 绘制 softmax 概率 \newline \#4 绘制 log-softmax 值}\par}

得到的绘图如图 5.14 所示。可以看到，对数概率与 logits、概率值的顺序相同。较大的 logit 对应较大的概率，以及不那么负的对数概率。反之，较小的 logit 产生较小的概率和更负的对数概率。

\myGraphic{0.92}{images/CH05_F14_Raschka2.png}{图 5.14 一个简单示例中 logits、softmax 概率和对数概率的比较。对数概率保留了概率的顺序。}

从实用角度来说，概率越高越好，因为它们表明模型对该词元赋予了更高的置信度。就对数概率而言，越接近零的值越好，因为零对应概率 1，而非常负的值对应极小的概率。这样就可以很方便地比较词元：最接近零的对数概率是概率最大者，最负的对数概率则是概率最小者。

图 5.14 用一个简单示例说明了 logits、概率和对数概率之间的关系，而我们可以把它应用到前面 \texttt{"}\texttt{The} \texttt{capital} \texttt{of} \texttt{Germany} \texttt{is}\texttt{"} 这个提示的概率值上，如图 5.15 所示。

\myGraphic{0.92}{images/CH05_F15_Raschka2.png}{图 5.15 在下一个词元评分中，logits 如何转换为概率和对数概率。正确的下一个词元（\texttt{"}\texttt{Berlin}\texttt{"}）获得较高的 logit，从而变成一个高概率和一个不那么负的对数概率，而不太可能的候选如 \texttt{"}\texttt{Bridge}\texttt{"}则对应极小的概率和绝对值很大的负对数概率。}

我们使用概率而不是 logits，因为概率是归一化的，更容易被理解为置信度值。前面我们还提到，对数概率还有一个额外的好处：与直接使用原始概率相比，它能带来数值上更稳定的计算。

数值稳定性来源于这样一个事实：即便是极小的概率分数，其对应的对数概率也落在一个合理的数值范围内。另一个原因是，把许多概率相乘会很快使结果趋近于零，而取对数能把这种乘法变成加法，从而避免下溢，对较长的序列也稳定得多。

图 5.16 展示了如何为我们前面用过的两段示例文本计算联合对数概率（或序列对数似然）。当我们计算普通联合概率分数时，看到两个序列的值都接近于零。现在，借助联合对数概率，我们可以看到，即使只是改变最后一个词元（例如 \texttt{"}\texttt{Berlin}\texttt{"} 与 \texttt{"}\texttt{Bridge}\texttt{"}），也会使总对数概率明显不同（\texttt{-16.6250} 与 \texttt{-29.8750}）。

\myGraphic{0.92}{images/CH05_F16_Raschka2.png}{图 5.16 词元级对数概率如何累加形成序列对数概率。每一行显示在给定前面文本的条件下下一个词元的对数概率。把这些值相加，就得到整个序列的联合对数概率。}

\begin{mySidebar}{处理对数概率}

当我们计算一个序列的联合概率时，我们会把许多介于 0 和 1 之间的数相乘：

\end{mySidebar}

\myGraphic{0.7}{images/eq-chapter-5-svg-eq-164.png}{}

\begin{mySidebarBox}

展开形式可以写成：

\end{mySidebarBox}

\myGraphic{0.7}{images/eq-chapter-5-svg-eq-166.png}{}

\begin{mySidebarBox}

这些乘积很快就会变得极其小，既不便于使用，还可能造成数值下溢。当数值变得过小时，计算机可能会把它向下舍入为零，我们就丢失了有用的信息。为避免这种情况，我们通常在 log 空间中进行处理。对联合概率取对数，就得到：

\end{mySidebarBox}

\myGraphic{0.7}{images/eq-chapter-5-svg-eq-168.png}{}

\begin{mySidebarBox}

利用“乘积的对数等于对数之和”这一性质，它可以展开为

\end{mySidebarBox}

\myGraphic{0.7}{images/eq-chapter-5-svg-eq-170.png}{}

\begin{mySidebarBox}

换句话说，一个序列的对数概率就是其各个词元对数概率之和。这既在数值上更稳定，也更便于处理。它还与 \texttt{torch.log\_softmax} 返回的值一致，后者直接给出对数概率。

因此，对对数概率求和是机器学习和 AI 中的标准做法。

\end{mySidebarBox}

实现词元对数概率函数很直接，因为只需对代码清单 5.6 中的代码做两处小的改动：把 \texttt{torch.softmax} 改为 \texttt{torch.log\_softmax}，并把 \texttt{torch.prod} 改为 \texttt{torch.sum}。

\begin{myWarning}{警告}
 使用不匹配的组合，例如 \texttt{torch.softmax} 配 \texttt{torch.sum}，或 \texttt{torch.log\_softmax} 配 \texttt{torch.prod}，技术上可行，但在数学上是错误的。\end{myWarning}

\begin{codeboxt}{代码清单 5.7 计算下一个词元的对数概率}
@torch.inference_mode()
def calc_next_token_logprobas(model, tokenizer, prompt, device, show=True):

    token_ids = torch.tensor(tokenizer.encode(prompt), device=device)

    logits = model(token_ids.unsqueeze(0)).squeeze(0)
    all_logprobas = torch.log_softmax(logits, dim=-1)    #1

    t_idx = torch.arange(0, token_ids.shape[0] - 1, device=device)
    next_ids = token_ids[1:]
    next_token_logprobas = all_logprobas[t_idx, next_ids]

    sum_next_token_logprobas = torch.sum(next_token_logprobas)    #2

    if show:
        print("Next-token log-probabilities:", next_token_logprobas)
        print("Joint log-probability:", sum_next_token_logprobas)
    else:
        return next_token_logprobas, sum_next_token_logprobas
\end{codeboxt}
{\noindent\small\textit{\#1 现在我们使用 log\_softmax。 \newline \#2 我们把乘积替换为求和。}\par}

现在，我们在前面的示例提示上试试这个更新后的函数：

\begin{codebox}
calc_next_token_logprobas(
        model, tokenizer, device=device,
        prompt="The capital of Germany is Berlin"
    )
\end{codebox}

结果如下：

\begin{codebox}
Next-token log-probabilities: tensor([-9.6875, -0.7695, -4.0938,
-0.3008, -1.7812], dtype=torch.bfloat16)
Joint log-probability: tensor(-16.6250, dtype=torch.bfloat16)
\end{codebox}

现在我们试试第二个序列：

\begin{codebox}
calc_next_token_logprobas(
        model, tokenizer, device=device,
        prompt="The capital of Germany is Bridge"
    )
\end{codebox}

返回如下结果：

\begin{codebox}
Next-token log-probabilities: tensor([ -9.6875,  -0.7695,  -4.0938,
-0.3008, -15.0000], dtype=torch.bfloat16)
Joint log-probability: tensor(-29.8750, dtype=torch.bfloat16)
\end{codebox}

如你所见，\texttt{"}\texttt{Berlin}\texttt{"} 与 \texttt{"}\texttt{Bridge}\texttt{"} 之间的差距现在明显得多（\texttt{-16.6250} 与 \texttt{-29.8750}），其中 \texttt{"}\texttt{Berlin}\texttt{"} 的得分高得多（更好）。

\section{用对数概率对模型置信度评分}

前两节非常详细地讲解了词元概率和对数概率的概念。本节将对对数概率的计算稍作修改，开发一个基于对数概率的评分函数，与本章开头开发的启发式评分器类似（见图 5.17）。

\myGraphic{0.92}{images/CH05_F17_Raschka2.png}{图 5.17 本节实现一个基于词元对数概率的 logprob 评分器，我们将在本章后面的自我修正方法中使用它。}

我们将要开发的 logprob 评分方法与上一节讲过的过程非常相似，但会做两处主要修改。第一，把提示排除在分数计算之外，只对答案词元计算分数。第二，对词元对数概率取平均（而不是求和），这样比较两个长度不同的序列会更公平。包含这两处修改的更新后计算如图 5.18 所示。

\myGraphic{0.92}{images/CH05_F18_Raschka2.png}{图 5.18 修改后的 logprob 评分过程。提示词元被排除在计算之外，只收集答案词元的对数概率。然后对这些值取平均，得到一个按长度归一化的分数，使我们能够公平地比较不同长度的答案。}

为便于说明，我们用代码清单 5.7 中的 \texttt{calc\_next\_token\_logprobas} 函数来实现图 5.18 所示的两处修改。例如，假设我们有下面这个提示和答案：

\begin{codebox}
example_prompt = "What is the capital of Germany?"
example_answer = " The capital of Germany is Berlin."

next_token_logprobas, sum_next_token_logprobas = calc_next_token_logprobas(
        model, tokenizer, device=device,
        prompt=example_prompt+example_answer,
        show=False
    )

print("Next-token logprobas:", next_token_logprobas)
print("Joint log-probability:", sum_next_token_logprobas)
\end{codebox}

它打印出如下输出。（注意，该张量也包含提示词元的分数，这就是数字看起来与图 5.18 不同的原因。继续读下去就会更清楚。）

\begin{codebox}
Next-token logprobas: tensor([-0.4512, -0.3418, -8.3125, -0.3906,
-3.8125, -3.0469, -1.1719,  0.0000, -0.0155,  0.0000,
-0.0078, -0.0752, -0.1582], dtype=torch.bfloat16)
Joint log-probability: tensor(-17.7500, dtype=torch.bfloat16)
\end{codebox}

接着，我们可以用下面的代码计算答案词元的数量，结果为 \texttt{7}：

\begin{codebox}
print(len(tokenizer.encode(example_answer)))
\end{codebox}

然后，为了按图 5.18 所示计算答案词元上的平均对数概率，我们需要对这七个答案词元取平均：

\begin{codebox}
last_7 = next_token_logprobas[-7:]
print(last_7)
print(torch.mean(last_7))
This prints:
tensor([-1.1719,  0.0000, -0.0155,  0.0000, -0.0078, -0.0752, -0.1582],
       dtype=torch.bfloat16)
tensor(-0.2041, dtype=torch.bfloat16)
\end{codebox}

可以看到，得到的平均值 \texttt{-0.2041} 与图 5.18 中所示相近。

我们可以把 \texttt{calc\_next\_token\_logprobas} 的代码与这一计算合并成一个新函数 \texttt{avg\_logprob\_answer}。\texttt{}

\begin{codeboxt}{代码清单 5.8 对答案词元的平均对数概率评分}
@torch.inference_mode()
def avg_logprob_answer(model, tokenizer, prompt, answer, device="cpu"):

    prompt_ids = tokenizer.encode(prompt)  #1
    answer_ids = tokenizer.encode(answer)   #2
    full_ids = torch.tensor(prompt_ids + answer_ids, device=device)

    logits = model(full_ids.unsqueeze(0)).squeeze(0)  #2
    logprobs = torch.log_softmax(logits, dim=-1)       #2

    start = len(prompt_ids) - 1  #3
    end = full_ids.shape[0] - 1   #3

    t_idx = torch.arange(start, end, device=device)#4
    next_tokens = full_ids[start + 1 : end + 1] #4
    next_token_logps = logprobs[t_idx, next_tokens] #4

    return torch.mean(next_token_logps)#5
\end{codeboxt}
{\noindent\small\textit{\#1 分别对提示和答案词元进行编码，以便稍后获取提示长度。 \newline \#2 与前面的 calc\_next\_token\_logprobas 相同 \newline \#3 与答案词元对应位置的索引范围 \newline \#4 除使用 start 和 end 外，其余与之前相同 \newline \#5 对答案词元的分数取平均。}\par}

\texttt{avg\_logprob\_answer} 函数与代码清单 5.7 中的 \texttt{calc\_next\_token\_logprobas} 函数相似，区别在于我们只对答案词元计算 logprob，而不是对包括提示在内的整个序列计算；并且对计算出的 logprob 值取平均而不是求和。

我们把这个新函数应用到本节前面的提示和答案上：

\begin{codebox}
score_1 = avg_logprob_answer(
    model, tokenizer,
    prompt="What is the capital of Germany?",
    answer=" The capital of Germany is Berlin.",
    device=device
)
print(score_1)
\end{codebox}

它和之前一样返回 \texttt{-0.2041}，说明我们的函数实现正确。现在我们也可以把这个函数应用到那个毫无意义的 \texttt{"}\texttt{Bridge}\texttt{"} 答案上：

\begin{codebox}
score_2 = avg_logprob_answer(
    model, tokenizer,
    prompt="What is the capital of Germany?",
    answer=" The capital of Germany is Bridge.",
    device=device
)
print(score_2)
\end{codebox}

这里得到的分数是 \texttt{-3.8906}，远低于 \texttt{"}\texttt{Berlin}\texttt{"} 的分数，符合预期。

有了这个平均 logprob 评分函数，我们也可以为本章开头定义的 MATH-500 提示（\texttt{prompt\_cot}）以及我们存为 \texttt{response\_1} 和 \texttt{response\_2} 的回复计算分数。这个就留给你作为练习了。

本章花了大量篇幅讲解下一个词元概率和对数概率评分的概念。这样做的一个原因是，我们可以在接下来关于自我修正的一节中把它们用作评分器。第二个原因是，对数概率这一概念在下一章实现带可验证奖励的强化学习训练时也会派上用场。

与此同时，重要的是不要以为 logprob 评分能包治百病。由于它衡量的是模型自身强烈偏好的东西，logprob 评分仍可能偏向那些自信的错误，而不选正确但偏好程度较低的答案。在实践中，最好把 logprob 评分看作若干有用评分信号中的一种，而不是可靠的平局裁决手段。

\begin{mySidebar}{练习 5.3：把 logprob 评分器用作自洽性中的平局裁决器}

仿照练习 5.1，但用 logprob 评分器替代启发式评分器，扩展上一章的自洽性实现（\texttt{self\_consistency\_vote}，见代码清单 4.17），使其能够处理候选答案之间的平局。然后在 MATH-500 数据集的一个子集上运行这两种实现，看看哪种平局裁决方法表现更好。

\end{mySidebar}

\begin{mySidebar}{练习 5.4：在 best-of-\emph{N} 设置中使用 logprob 评分器}

仿照练习 5.2，扩展自洽性实现，使最终答案用 logprob 评分器（\texttt{avg\_logprob\_answer}）来选出，而不是用启发式评分器或多数投票。然后在 MATH-500 数据集的一个子集上运行这些不同的实现，看看哪种平局裁决方法表现更好。

提示：无需修改 \texttt{self\_consistency\_vote }函数本身来实现平局裁决。你可以对 \texttt{self\_consistency\_vote} 返回的 \texttt{results} 字典进行平局裁决，该字典在代码清单 3.15 的 \texttt{evaluate\_math500\_stream} 函数内部被使用。

\end{mySidebar}

\section{通过迭代反馈实现自我修正}

在介绍了几种为 LLM 答案评分的方法之后，现在来看看本章的核心推理扩展技术——\emph{自我修正}（图 5.19）。我们先介绍自我修正，并手动走一遍整个流程。在下一节，如图 5.19 所示，我们会把自我修正循环自动化，并为本章前面开发的评分方法（启发式评分器和平均 logprob 评分器）添加支持。

\myGraphic{0.92}{images/CH05_F19_Raschka2.png}{图 5.19 我们工作流程中的最后一步，其中前面开发的 logprob 评分器被用于自我修正方法}

自我修正是一种让 LLM 分析和改进自己答案的技术。如图 5.20 所示，LLM 先像常规使用那样给出对提示的初始答案，然后对答案进行批判并加以改进。

图 5.20 中的自我修正流程乍看可能有些复杂，但它本质上只是把文本生成函数依次应用到不同的提示和输入上。我们用一个代码示例一步步走一遍。

\begin{myWarning}{注意}
 在本例中，我们不会使用思维链提示。不过在实践中，当使用基础模型时，你可以把自我修正与思维链提示结合起来。\end{myWarning}

我们将从基础提示和答案开始（图 5.20 中的步骤 1 和步骤 2），基于 5.2 节的代码。

\myGraphic{0.92}{images/CH05_F20_Raschka2.png}{图 5.20 自我修正流程。LLM 先对提示给出一个初始答案，然后收到一个批判提示（critique prompt），要求它分析自己的回复，并给出一个简短的批判以及改进答案的计划。在最后一步，模型会收到一个改进提示（refine prompt），其中包含原始问题、它的草稿答案以及批判内容，模型据此生成一个纳入建议改进的修订答案。}

\begin{codeboxt}{代码清单 5.9 基础提示和答案}
raw_prompt = (
    "Half the value of $3x-9$ is $x+37$. "
    "What is the value of $x$?"
)
prompt = render_prompt(raw_prompt)

torch.manual_seed(123)
initial_response = generate_text_stream_concat_flex(
    model, tokenizer, prompt, device,
    max_new_tokens=2048, verbose=True,
    generate_func=generate_text_top_p_stream_cache,
    temperature=0.7,
    top_p=0.9,
)
The LLM responds with " \boxed{18}", which is incorrect (the correct answer is 83).
\end{codeboxt}

接下来，我们让 LLM 对答案进行批判。为此，我们编写一个批判提示，其中包含原始问题（\texttt{raw\_prompt}）和答案（\texttt{draft}），如图 5.20 中的步骤 3 和步骤 4 所示。

\begin{codeboxt}{代码清单 5.10 批判提示和改进计划}
def make_critique_prompt(raw_prompt, draft):
    return (
        "You are a meticulous reviewer. Identify logical errors, missing "
        "steps, or arithmetic mistakes. If the answer seems correct, "
        "say so briefly. Then propose a concise plan to fix issues.\n\n"
        f"Question:\n{raw_prompt}\n\n"
        f"Draft answer:\n{draft}\n\n"
        "Write a short critique and bullet-point fix plan "
        "(under ~120 words).\n"
        "Critique:"
    )

critique_prompt = make_critique_prompt(raw_prompt, initial_response)
torch.manual_seed(123)
critique = generate_text_stream_concat_flex(
    model, tokenizer, critique_prompt, device,
    max_new_tokens=2048, verbose=True,
    generate_func=generate_text_top_p_stream_cache,
    temperature=0.7,
    top_p=0.9,
)
\end{codeboxt}

上述批判提示让 LLM 写出一份相当详尽的批判，其中甚至包含正确答案本身：

\begin{codebox}
 The question seems to have a logical error in its setup. The statement "Half the value of $3x-9$ is $x+37$" is incorrect because half of $3x-9$ should be $(3x-9)/2$, not $3x-9$. The equation should be $\frac{1}{2}(3x-9) = x + 37$.

Fix Plan:
1. Correct the equation to $\frac{1}{2}(3x-9) = x + 37$.
2. Multiply both sides by 2 to eliminate the fraction: $3x - 9 = 2(x + 37)$.
3. Distribute the 2 on the right side: $3x - 9 = 2x + 74$.
4. Subtract $2x$ from both sides: $x - 9 = 74$.
5. Add 9 to both sides: $x = 83$.
\end{codebox}

注意，批判中可能包含事实性错误。例如，它说 \texttt{"}\texttt{The} \texttt{question itself} \texttt{is} \texttt{incomplete}\texttt{"}，这并不属实，但它仍然接着给出了正确的修复计划。这一点很重要，因为在自我修正中，即便批判的一部分有误，只要总体上能推动模型做出更好的修订，它就是有用的。

最后，我们用这个改进提示来修订原始答案（图 5.20 中的步骤 5 和步骤 6）。

\begin{codeboxt}{代码清单 5.11 答案改进}
def make_refine_prompt(raw_prompt, draft, critique):
    return (
        "Revise the answer using the critique. Keep it concise and "
        "end with a final boxed result: \\boxed{ANSWER}\n\n"
        f"Question:\n{raw_prompt}\n\n"
        f"Previous answer:\n{draft}\n\n"
        f"Critique:\n{critique}\n\n"
        "Revised answer:"
    )

refine_prompt = make_refine_prompt(raw_prompt, initial_response, critique)
torch.manual_seed(123)
revised_answer = generate_text_stream_concat_flex(
    model, tokenizer, refine_prompt, device,
    max_new_tokens=2048, verbose=True,
    generate_func=generate_text_top_p_stream_cache,
    temperature=0.7,
    top_p=0.9,
)
\end{codeboxt}

虽然基础模型并不是最好的指令遵循者，但改进提示仍然让 LLM 生成了正确的答案：

\begin{codebox}
...
 \boxed{83}
Final result: The value of $x$ is \boxed{83}.
\end{codebox}

这个示例演示了单次改进循环。在实践中，通常会把这个循环重复多次迭代。接下来，我们将编写一个函数来自动化这一过程。

\section{编写自我修正循环}

现在，我们把上一节手动执行的自我修正步骤打包成一个方便使用的函数，从而简化自我修正的使用，并让这一过程能够按固定的迭代次数重复进行。

此外，该函数还将支持接入本章前面开发的评分方法（启发式评分和 logprob 评分），为每个答案计算分数（图 5.21）。这个分数随后可用来决定是否应当采纳某个改进后的答案。

\myGraphic{0.92}{images/CH05_F21_Raschka2.png}{图 5.21 带可选评分的自我修正循环。模型给出一个初始答案，然后对它进行批判，并基于批判生成一个修订答案。两个答案都可以用评分函数（如 logprob 评分）来评估，只有当修订答案的分数优于前一个答案时，才予以采纳。}

在图 5.21 中，修订答案得到的 logprob 分数（\texttt{-0.377}）高于初始答案（\texttt{-1.258}），因此采纳这一修订是合理的。在实践中，自我修正也可能产生更差的答案，尤其是在运行多次迭代时。因此，评分器提供了一种判断修订答案是否带来改进的方法。

\begin{myWarning}{注意}
 评分并不总能改善结果。在自我修正中是否应当使用评分器、以及该用哪一个，取决于具体的 LLM。你可以通过在 MATH-500 这样的基准数据集上做实验来判断。\end{myWarning}

把上一节的各个步骤合并起来、并加入 \texttt{iterations} 参数和评分函数（\texttt{score\_fn}）选项的完整代码，见下面的代码清单。

\begin{codeboxt}{代码清单 5.12 支持多次迭代和评分的自我修正}
def self_refinement_loop(
    model,
    tokenizer,
    raw_prompt,
    device,
    iterations=2,
    max_response_tokens=2048,
    max_critique_tokens=256,
    score_fn=None,
    prompt_renderer=render_prompt,
    prompt_suffix="",
    verbose=False,
    temperature=0.7,
    top_p=0.9,
):
    steps = []

    prompt = prompt_renderer(raw_prompt) + prompt_suffix #1
    current_full = generate_text_stream_concat_flex(
        model=model,
        tokenizer=tokenizer,
        prompt=prompt,
        device=device,
        max_new_tokens=max_response_tokens,
        verbose=False,
        generate_func=generate_text_top_p_stream_cache,
        temperature=temperature,
        top_p=top_p,
    )

    current_extracted = extract_final_candidate(
        current_full, fallback="number_then_full"
    )
    if score_fn:
        current_score = score_fn(answer=current_full, prompt=prompt)
    else:
        current_score = 0.0

    for it in range(iterations):     #2
        draft_before_full = current_full
        draft_before_extracted = current_extracted
        score_before = current_score

        critique_prompt = make_critique_prompt(        #3
            raw_prompt, draft_before_full
        )
        critique_full = generate_text_stream_concat_flex(
            model=model,
            tokenizer=tokenizer,
            prompt=critique_prompt,
            device=device,
            max_new_tokens=max_critique_tokens,
            verbose=False,
            generate_func=generate_text_top_p_stream_cache,
            temperature=temperature,
            top_p=top_p,
        )

        refine_prompt = make_refine_prompt(         #4
            raw_prompt, draft_before_full, critique_full
        )
        revised_full = generate_text_stream_concat_flex(
            model=model,
            tokenizer=tokenizer,
            prompt=refine_prompt,
            device=device,
            max_new_tokens=max_response_tokens,
            verbose=False,
            generate_func=generate_text_top_p_stream_cache,
            temperature=temperature,
            top_p=top_p,
        )

        revised_extracted = extract_final_candidate(
            revised_full, fallback="number_then_full"
        )
        if score_fn:
            revised_score = score_fn(
                answer=revised_full, prompt=prompt
            )
        else:
            revised_score = 0.0

        step = {         #5
            "iteration": it + 1,
            "draft_full": draft_before_full,
            "draft_extracted": draft_before_extracted,
            "critique": critique_full,
            "revised_full": revised_full,
            "revised_extracted": revised_extracted,
            "score_before": score_before,
            "score_after": revised_score,
        }
        steps.append(step)

        if verbose:
            print(
                f"[Refinement {it+1}/{iterations}]"
                f"\nCurrent: {draft_before_extracted}"
                f"\nRevised: {revised_extracted}"
                f"\nScore before: {score_before:.3f}"
                f"\nScore after: {revised_score:.3f}"
                f"\n{'=' * 25}\n"
            )

        if revised_score >= current_score:         #6
            current_full = revised_full
            current_extracted = revised_extracted
            current_score = revised_score

    return {
        "final_full": current_full,
        "final_extracted": current_extracted,
        "steps": steps,
    }
\end{codeboxt}
{\noindent\small\textit{\#1 初始回复（草稿） \newline \#2 运行一次或多次迭代。 \newline \#3 对回复进行批判。 \newline \#4 改进回复。 \newline \#5 记录结果。 \newline \#6 如果修订后的回复不比原来差，就采纳它。}\par}

\texttt{self\_refinement\_loop} 函数把上一节的自我修正流程运行一次或多次迭代（\texttt{for it in range(iterations)}）。在每次迭代结束时，它把修订分数与当前分数比较（\texttt{revised\_score \textgreater{}= current\_score}），只有当修订答案的分数等于或更高时才保留它。

评分函数的使用是可选的。当 \texttt{score\_fn=None}（默认值）时，分数始终设为 \texttt{0.0}。由于 \texttt{0.0 \textgreater{}= 0.0} 求值为 \texttt{True}，因此在运行多次改进迭代时，最近的答案总是会被采纳。

接着，我们用平均 logprob 评分器 \texttt{avg\_logprob\_answer} 来运行自我修正循环。这个函数（见代码清单 5.8）接受若干参数（\texttt{model}、\texttt{tokenizer}、\texttt{prompt}、\texttt{answer} 和 \texttt{device}）。如上一个代码清单所示，评分函数在调用时只传入了两个参数：

\begin{codebox}
# ...
        if score_fn:
            revised_score = score_fn(
                answer=revised_full, prompt=prompt
            )
# ...
\end{codebox}

为了让 \texttt{avg\_logprob\_answer} 与这种 \texttt{score\_fn} 调用方式兼容，我们可以使用 Python 内置 \texttt{functools} 模块中的 \texttt{partial} 函数来预先指定其余参数。

\begin{codeboxt}{代码清单 5.13 创建平均对数概率评分器}
from functools import partial

avg_logprob_score = partial(
    avg_logprob_answer,
    model=model,
    tokenizer=tokenizer,
    device=device
)
\end{codeboxt}

现在我们可以在自我修正循环中使用 \texttt{avg\_logprob\_score} 函数：

\begin{codebox}
torch.manual_seed(1)

results_logprob = self_refinement_loop(
    model=model,
    tokenizer=tokenizer,
    raw_prompt=raw_prompt,
    device=device,
    iterations=2,
    max_response_tokens=2048,
    max_critique_tokens=256,
    score_fn=avg_logprob_score,
    verbose=True,
    temperature=0.7,
    top_p=0.9,
)
\end{codebox}

输出如下：

\begin{codebox}
[Refinement 1/2]
Current: 10
Revised: 83
Score before: -0.855
Score after: -0.226
=========================
[Refinement 2/2]
Current: 83
Revised: 83
Score before: -0.226
Score after: -1.320
=========================
\end{codebox}

观察上面的结果可以看到，模型在第一次迭代中纠正了最初错误的答案（从 \texttt{10} 改为 \texttt{83}），分数从 \texttt{-0.855} 提升到 \texttt{-0.226}。在第二次迭代中，分数变差了，不过没关系，因为我们已经有了正确的答案。

我们还可以通过 \texttt{results\_logprob[}\texttt{"}\texttt{final\_extracted}\texttt{"}\texttt{]} 访问从得分最高的回复中提取出的最终答案（通常来自其 boxed 内容，如果有的话），在本例中它返回 \texttt{83}。如果你对完整答案感兴趣，请使用 \texttt{results\_logprob[}\texttt{"}\texttt{final\_full}\texttt{"}\texttt{]}。要浏览批判和更长的答案，可以打印 \texttt{results\_logprob} 字典，其中包含所有迭代的详细结果。

\begin{mySidebar}{练习 5.5：用启发式评分进行自我修正}

使用我们在代码清单 5.3 中定义的 \texttt{heuristic\_score} 函数来运行 \texttt{self\_refinement\_loop}。

\end{mySidebar}

前面的示例表明，自我修正可以帮助模型生成正确答案。为了更清楚地了解这种自我修正方法究竟有多大用处，我在 MATH-500 上运行了它，并把结果汇总在表 5.1 中。

\begin{table}[htbp]
\centering\small
\centering
\caption{表 5.1 不同自我修正方法在 MATH-500 任务上的准确率}
\begin{tabularx}{\linewidth}{@{}c@{}XXXXXX}
\toprule
 & 方法 & 评分 & 迭代次数 & 模型 & 准确率 & 时间 \\
1 & 基线（第 3 章） & - & - & Base & 15.2\% & 10.1 min \\
2 & 自我修正 & None & 1 & Base & 25.0\% & 84.8 min \\
3 & 自我修正 & None & 2 & Base & 22.0\% & 165.4 min \\
4 & 自我修正 & 启发式 & 1 & Base & 21.6\% & 84.7 min \\
5 & 自我修正 & 启发式 & 2 & 基础模型 & 20.8\% & 151.4 min \\
6 & 自我修正 & 平均 logprob & 1 & 基础模型 & 21.4\% & 85.3 min \\
7 & 自我修正 & 平均 logprob & 2 & 基础模型 & 22.0\% & 165.3 min \\ \addlinespace

8 & 基线（第 3 章） & - & - & 推理模型 & 48.2\% & 182.1 min \\
9 & 自我修正 & None & 1 & 推理模型 & 56.6\% & 498.8 min \\
10 & 自我修正 & 启发式 & 1 & 推理模型 & 57.8\% & 498.6 min \\
11 & 自我修正 & 平均 logprob & 1 & 推理模型 & 48.4\% & 499.7 min \\ \bottomrule
\end{tabularx}
\end{table}

表 5.1 中所示的准确率值是在 MATH-500 测试集全部 500 个样本上、使用 \texttt{"}\texttt{cuda}\texttt{"} GPU（DGX Spark）计算得到的。

如表 5.1 所示，自我修正在基础模型之上提升了性能（第 1–7 行），但提升幅度非常有限，而且最佳准确率出现在自我修正中不使用评分时。这意味着启发式评分和平均 logprob 评分有时会导致错误的答案被采纳，取代了最初正确的答案。

注意，当向基础模型和自我修正都加上 \texttt{"}\texttt{Explain} \texttt{step} \texttt{by} \texttt{step.}\texttt{"} 思维链时，模型未能比基础模型有所提升（结果未展示）。

再看推理模型的结果（第 8–11 行），可以看到自我修正与启发式评分的组合把答案准确率提升了近 10\%。（如第 3 章所解释的，要使用推理模型，只需在代码清单 5.1 的 \texttt{load\_model\_and\_tokenizer} 中把 \texttt{which\_model=}\texttt{"}\texttt{base}\texttt{"} 改为 \texttt{which\_model=}\texttt{"}\texttt{reasoning}\texttt{"}。）

在这两种情况下，平均 logprob 评分似乎都比不用评分器或使用启发式评分器表现更差。这很可能是因为，平均 logprob 分数与一个答案在模型看来有多自然或多符合预期关系更密切，而与答案是否真的正确关系较小。因此，它可能偏好流畅、熟悉或句法整洁的答案，即便这些答案在语义上是错误的。换句话说，logprob 准则可能无意中选中那些自信的错误，而启发式评分则更关注答案的格式和结构。

到这里，可能看起来我们在讨论平均 logprob 评分这一概念上花了太多时间。但在处理 LLM 时，logprob 评分是一个基础概念，在接下来的章节中实现强化学习训练流程时，它还会派上用场。

把表 5.1 的结果与上一章的结果（表 4.1）相比，我们还观察到，在这个数学任务上，对该模型而言自我修正不如自洽性有效。自洽性仍然被广泛使用，因为它尽管简单，却在实践中表现良好。

本章尚未探讨的另一种方法是，在 LLM 作为评判者的设置中用外部模型进行自我修正，与附录 F（F.5 节）中所述类似。例如，我们可以不用启发式评分或平均 logprob，而是用第二个 LLM 来计算分数并撰写批判。

2025 年 11 月，DeepSeek 团队借助 DeepSeekMath-V2 证明，用第二个 LLM 进行自我修正可以非常成功，在多项数学竞赛中达到金牌水平。具体来说，该团队提出了一种新方法：他们训练第二个 LLM 成为一个强大的批判模型，并进一步训练其基础模型，使其在自我修正的场景下成为更好的数学解题者（相比之下，本章中的自我修正是完全不做额外训练的）。

至此，我们对无需额外训练的推理扩展的讨论就结束了（见图 5.22）。下一章，我们将开始介绍通过更新模型权重来提升推理的各类训练技术。

\myGraphic{0.92}{images/CH05_F22_Raschka2.png}{图 5.22 本书从基础 LLM 使用，到推理时推理方法，再到基于训练的技术这一演进脉络的概览。本章结束了我们对无需额外训练的推理扩展方法的讨论。接下来的几章将介绍通过更新模型权重来进一步提升推理性能的方法。}

\section{小结}

\begin{itemize}
\item 自我修正延续了上一章的推理时扩展思路，通过对单个答案反复批判和改进来提升效果，而不是像自洽性那样依赖多个独立样本。
\item 一个简单的基于规则的评分函数会对模型输出排序，它对可提取的最终答案以及更短、更经济的补全给予奖励。
\item 下一个词元评分把 logits 转换为归一化概率，并将它们合并为序列级似然，从而量化模型的置信度。
\item 对数概率取代原始概率，以避免数值下溢，并把跨多个词元的乘积转变为稳定的求和与求平均。
\item 这些评分函数的用途不限于自我修正，例如还可用于在自洽性中打破平局，或实现 best-of-\emph{N} 选择策略。
\item 自我修正流程包含三个阶段：生成初始草稿、给出简短的批判和修复计划，以及生成修订后的答案。
\item 一个可复用的改进函数把自我修正工作流自动化，跨越多次迭代（改进轮次）；它还可以借助基于分数的接纳机制，只保留那些不会让所计算答案分数下降的修订——分数由我们前面开发的某个评分函数计算得出。
\end{itemize}
